% problem_id: S001-theorem-zariski-fpqc-quasi-coherent
% source_id: S001 (Stacks Project)
% source_file: etale-cohomology.tex
% statement_locator: etale-cohomology.tex lines 2279--2297
% proof_locator: etale-cohomology.tex lines 2299--2375
% env: theorem
% id_note: label 在本来源内唯一
% status: complete (statement + author proof verbatim + reason + locators)
%
\begin{theorem}
\label{theorem-zariski-fpqc-quasi-coherent}
Let $S$ be a scheme and $\mathcal{F}$ a quasi-coherent $\mathcal{O}_S$-module.
Let $\mathcal{C}$ be either $(\Sch/S)_\tau$ for
$\tau \in \{fppf, syntomic, smooth, \etale, Zariski\}$ or
$S_\etale$. Then
$$
H^p(S, \mathcal{F}) = H^p_\tau(S, \mathcal{F}^a)
$$
for all $p \geq 0$ where
\begin{enumerate}
\item the left hand side indicates the usual cohomology of the sheaf
$\mathcal{F}$ on the underlying topological space of the scheme $S$, and
\item the right hand side indicates cohomology
of the abelian sheaf $\mathcal{F}^a$ (see
Proposition \ref{proposition-quasi-coherent-sheaf-fpqc})
on the site $\mathcal{C}$.
\end{enumerate}
\end{theorem}

\begin{proof}
We are going to show that
$H^p(U, f^*\mathcal{F}) = H^p_\tau(U, \mathcal{F}^a)$
for any object $f : U \to S$ of the site $\mathcal{C}$.
The result is true for $p = 0$ by the sheaf property.

\medskip\noindent
Assume that $U$ is affine. Then we want to prove that
$H^p_\tau(U, \mathcal{F}^a) = 0$ for all $p > 0$. We use induction on $p$.
\begin{enumerate}
\item[p = 1]
Pick $\xi \in H^1_\tau(U, \mathcal{F}^a)$.
By Lemma \ref{lemma-locality-cohomology},
there exists an fpqc covering $\mathcal{U} = \{U_i \to U\}_{i \in I}$
such that $\xi|_{U_i} = 0$ for all $i \in I$. Up to refining
$\mathcal{U}$, we may assume that $\mathcal{U}$ is a standard
$\tau$-covering. Applying the spectral sequence of
Theorem \ref{theorem-cech-ss},
we see that $\xi$ comes from a cohomology class
$\check \xi \in \check H^1(\mathcal{U}, \mathcal{F}^a)$.
Consider the covering $\mathcal{V} = \{\coprod_{i\in I} U_i \to U\}$. By
Lemma \ref{lemma-cech-complex},
$\check H^\bullet(\mathcal{U}, \mathcal{F}^a) =
\check H^\bullet(\mathcal{V}, \mathcal{F}^a)$.
On the other hand, since $\mathcal{V}$ is a covering of the form
$\{\Spec(B) \to \Spec(A)\}$ and $f^*\mathcal{F} = \widetilde{M}$
for some $A$-module $M$, we see the {\v C}ech complex
$\check{\mathcal{C}}^\bullet(\mathcal{V}, \mathcal{F})$
is none other than the complex $(B/A)_\bullet \otimes_A M$.
Now by Lemma \ref{lemma-descent-modules},
$H^p((B/A)_\bullet \otimes_A M) = 0$ for $p > 0$, hence $\check \xi = 0$
and so $\xi = 0$.
\item[p > 1]
Pick $\xi \in H^p_\tau(U, \mathcal{F}^a)$. By
Lemma \ref{lemma-locality-cohomology},
there exists an fpqc covering $\mathcal{U} = \{U_i \to U\}_{i \in I}$
such that $\xi|_{U_i} = 0$ for all $i \in I$. Up to refining
$\mathcal{U}$, we may assume that $\mathcal{U}$ is a standard
$\tau$-covering. We apply the spectral sequence of
Theorem \ref{theorem-cech-ss}.
Observe that the intersections $U_{i_0} \times_U \ldots \times_U U_{i_p}$
are affine, so that by induction hypothesis the cohomology groups
$$
E_2^{p, q} = \check H^p(\mathcal{U}, \underline{H}^q(\mathcal{F}^a))
$$
vanish for all $0 < q < p$. We see that $\xi$ must come from a
$\check \xi \in \check H^p(\mathcal{U}, \mathcal{F}^a)$. Replacing
$\mathcal{U}$ with the covering $\mathcal{V}$ containing only one morphism
and using Lemma \ref{lemma-descent-modules} again,
we see that the {\v C}ech cohomology class $\check \xi$ must be zero,
hence $\xi = 0$.
\end{enumerate}
Next, assume that $U$ is separated. Choose an affine open covering
$U = \bigcup_{i \in I} U_i$ of $U$. The family
$\mathcal{U} = \{U_i \to U\}_{i \in I}$ is then an fpqc covering,
and all the intersections
$U_{i_0} \times_U \ldots \times_U U_{i_p}$ are affine
since $U$ is separated. So all rows of the spectral sequence of
Theorem \ref{theorem-cech-ss}
are zero, except the zeroth row. Therefore
$$
H^p_\tau(U, \mathcal{F}^a) =
\check H^p(\mathcal{U}, \mathcal{F}^a) =
\check H^p(\mathcal{U}, \mathcal{F}) = H^p(U, \mathcal{F})
$$
where the last equality results from standard scheme theory, see
Cohomology of Schemes, Lemma
\ref{coherent-lemma-cech-cohomology-quasi-coherent}.

\medskip\noindent
The general case is technical and (to extend the proof as given here)
requires a discussion about maps of spectral sequences, so we won't treat it.
It follows from
Descent, Proposition \ref{descent-proposition-same-cohomology-quasi-coherent}
(whose proof takes a slightly different approach) combined with
Cohomology on Sites, Lemma \ref{sites-cohomology-lemma-cohomology-of-open}.
\end{proof}
